\documentclass[11pt,a4paper]{article}
\usepackage[margin=1in]{geometry}
\usepackage{amsmath,booktabs,hyperref,float}
\definecolor{good}{rgb}{0.0,0.5,0.0}

\title{Iteration 009: Spatio-Temporal FIR on Real Data}
\author{Autoresearch Loop}
\date{April 7, 2026}

\begin{document}
\maketitle

\begin{abstract}
Adding temporal context via FIR filters improves scalp-to-in-ear prediction
from $r=0.366$ to $r=0.373$ (+1.9\%), the first improvement over the
closed-form baseline on real data. At 20\,Hz, a short FIR filter captures
50--550\,ms of temporal dynamics in the volume conduction path. Center-tap
initialization from the closed-form solution remains critical.
\end{abstract}

\section{Method}
Spatio-temporal FIR filter: $\hat{y}_c[t] = \sum_j \sum_\tau w_{c,j,\tau} x_j[t-\tau]$,
implemented as \texttt{nn.Conv1d}. Swept filter lengths $L \in \{3,5,7,11\}$ taps
(150--550\,ms at 20\,Hz). Center tap initialized from $\mathbf{W}^*$, all other
taps initialized to zero. Trained 150 epochs with AdamW ($\text{lr}=10^{-3}$),
cosine annealing, gradient clipping (norm 1.0).

\section{Results}
\begin{table}[H]
\centering
\caption{Benchmark results (subjects 13--15 LOSO).}
\begin{tabular}{lrrr}
\toprule
Model & Mean $r$ & Std $r$ & SNR (dB) \\
\midrule
Closed-form baseline & 0.366 & 0.072 & 0.59 \\
\textbf{FIR (best $L$)} & \textbf{0.373} & 0.074 & \textbf{0.61} \\
\bottomrule
\end{tabular}
\end{table}

\section{Analysis}
The +0.007 improvement is modest but consistent across test subjects, confirming
that temporal lags carry real information. At 20\,Hz with 1--9\,Hz content,
the temporal autocorrelation is strong and even a few taps capture useful
cross-lag structure.

The improvement ceiling at this sampling rate is likely low --- the 1--9\,Hz
band has limited spatial discriminability. Broadband data (1--45\,Hz) would
provide much richer temporal/spatial features.

\section{Next}
Iteration 010 will explore a deeper temporal model: a lightweight 1D
convolutional network with residual connections, larger receptive field,
and batch normalization. The hypothesis is that stacking temporal convolutions
can capture nonlinear temporal dynamics the linear FIR misses.

\end{document}
